\documentclass[addpoints]{exam}
% \documentclass[addpoints,12pt]{exam}

\usepackage[slovene]{babel}
\usepackage{amsfonts}
\usepackage{amssymb}
\usepackage{amsmath}
\usepackage{float}
\usepackage{graphicx}
\usepackage{enumitem}
\usepackage{color}
\usepackage[gen]{eurosym}
\usepackage{newunicodechar}
\newunicodechar{€}{\euro}


%Komentar naslednje vrstice glede na verzijo testa -- rešitve ali prostor za odgovore:
% \printanswers

% \linespread{2}


\pointsinrightmargin
\marginbonuspointname{}


\renewcommand{\solutiontitle}{\noindent\textbf{Rešitve in točkovanje:}\par\noindent}

\colorgrids
\definecolor{GridColor}{gray}{0.7}
\setlength{\gridsize}{7mm}

\extrawidth{5mm}
\setlength{\rightpointsmargin}{1.5cm}

\begin{document}

\runningfooter{}{}{}

\begin{center}
    \fbox{\fbox{\parbox{15cm}{\centering
    Popravljanje in pridobivanje ocen -- konferenca \\ 1. letnik -- test A \\ 12. junij 2025 \\ (Deljivost; Racionalna števila; Koordinatni sistem v ravnini; Funkcija; Premica)}}}
\end{center}

\vspace{0.1in}
\makebox[\textwidth]{Ime in priimek, oddelek:\enspace\hrulefill}


\begin{center}
    \hqword{Naloga}
    \hpword{Možne točke}
    \chbpword{Dodatne točke}
    \hsword{Dosežene točke}
    \htword{Skupaj}  
    % \combinedpointtable[h][questions]
    \gradetable[h][questions]

\end{center}

\vspace{0.1in}


\begin{questions}


    % 1. vprašanje

    \question[4]
    Določite največji skupni delitelj in najmanjši skupni večkratnik izrazov $x^3-4x^2+3x$ in $x(x^2-6x+9)$.

    \begin{solution}[7cm]
        Za pravilno faktorizacijo izrazov $x(x-2)(x+2)$ in $(x-2)(x^2+2x+4)$ po $1$ točka. 
        Za zapis največjega skupnega delietlja $x-2$ $1$ točka in najmanjšega skupnega večkratnika $x(x-2)(x+2)(x^2+2x+4)$ $1$ točka.
    \end{solution}
    

    

    % \newpage
    % 5. vprašanje
    \question[5]
    Izračunajte vrednost izraza.
        $$\left(0.\overline{15}\cdot 6\frac{3}{5}-2.2\right)\cdot\left(-0.8\overline{33}\right) $$

        \begin{solution}[4cm]
            \begin{align*}
                5.\overline{6}:(2.5-2.\overline{3}) - 0.3\overline{72}\cdot 110&=5\frac{2}{3}:(2\frac{1}{2}-2\frac{1}{3})- \frac{41}{110}\cdot 110 = \\
                                                                            &=\frac{17}{3}:\frac{1}{6}-41= \\
                                                                            &=\frac{17}{3}\cdot \frac{6}{1}-41= \\
                                                                            &=34-41= \\
                                                                            &= -7
            \end{align*}
            Za pravilno zapisano število $0.3\overline{72}$ z ulomkom $1$ točka, pravilno zapisani števili $5.\overline{6}$ in $1.\overline{3}$ z ulomkoma $1$ točka, pravilno množenje in deljenje ulomkov $1$ točka, pravilno odštevanje ulomkov $1$ točka, končen rezultat $1$ točka.
        \end{solution}




    \newpage
    % 6. vprašanje

    \question
    Okrajšajte ulomke.
    \begin{parts}
        

        \part[4]
        $\left(\dfrac{z+1}{z-1}+1+\dfrac{z}{z-1}\right)\cdot\left(\dfrac{3z+3}{z^2-1}\right)^{-1}$

        \begin{solution}[7cm]
            \begin{align*}
                \dfrac{x+2}{(x-2)^{-1}}:\left(\dfrac{1}{x^2-4}\right)^{-1}&= \dfrac{(x-2)(x+2)}{1}\cdot\dfrac{1}{x^2-4}= \\
                                                                            &= \dfrac{x^2-4}{x^2-4}= \\
                                                                            &= 1
            \end{align*}
            Za pravilno upoštavanje obratne vrednosti $1$ točka; pravilno deljenje ulomkov $1$ točka; pravilna faktorizacija oziroma razčlenjevanje $1$ točka; pravilen okrajšan rezultat $1$ točka.
        \end{solution}


        \part[5]
        $\dfrac{3^{1-2n}+5\cdot 3^{2-2n}+2\cdot 3^{3-2n}}{3^{1-3n}+3^{2-3n}}$

        \begin{solution}[6cm]
            \begin{align*}
                \dfrac{7^{3n-2}-7^{3n-1}}{7^{3n-3}+3\cdot 7^{3n-2}-7^{3n-4}} &= \dfrac{7^{3n-2}\left(1-7^1\right)}{7^{3n-4}\left(7^1+3\cdot 7^2-1\right)}= \\
                                                                            &= \dfrac{7^{3n-2}\left(1-7\right)}{7^{3n-4}\left(7+3\cdot 49-1\right)}= \\
                                                                            &= \dfrac{7^{3n-2}\cdot(-6)}{7^{3n-4}\cdot 153}= \\
                                                                            &= \dfrac{7^{3n-2}}{7^{3n-4}}\cdot\dfrac{-2}{51}= \\
                                                                            &= -\dfrac{2\cdot 7^2}{51}
            \end{align*}
            Za pravilno izpostavljanje skupnega faktorja v števcu in imenovalcu po $1$ točka, pravilno seštevanje $1$ točka, pravilno krajšanje potenc $1$ točka, zapis pravilno okrajšanega ulomka $1$ točka.
        \end{solution}



    \end{parts}




        % 7. vprašanje

        \question[5]
        Poenostavite izraz.
            $$1-x+x^2-\dfrac{1}{x^{-3}+x^{-2}}; \quad x\neq -1$$
    
            \begin{solution}[4cm]
                \begin{align*}
                    \dfrac{2x^{-2}}{1-2x^{-1}+x^{-2}}+\dfrac{x^{-1}}{1-x^{-1}} &= \dfrac{\frac{2}{x^2}}{1-\frac{2}{x}+\frac{1}{x^2}}+\dfrac{\frac{1}{x}}{1-\frac{1}{x}}= \\
                                                                            &= \dfrac{\frac{2}{x^2}}{\frac{x^2-2x+1}{x^2}}+\dfrac{\frac{1}{x}}{\frac{x-1}{x}}= \\
                                                                            &= \dfrac{2x^2}{(x^2-2x+1)x^2}+\dfrac{x}{(x-1)x}= \\
                                                                            &= \dfrac{2}{(x-1)^2}+\dfrac{1}{x-1}= \\
                                                                            &= \dfrac{2+(x-1)}{(x-1)^2}= \\
                                                                            &= \dfrac{x+1}{(x-1)^2}
                \end{align*}
                Za pravilen zapis obratnih vrednosti $1$ točka, pravilno seštevanje in odštevanje ulomkov $1$ točka, pravilno razreševanje dvojnih ulomkov $1$ točka, pravilno krajšanje ulomkov $1$ točka, končen rezultat $1$ točka.
            \end{solution}
    










\newpage


        % 3. vprašanje
    \question
    Dana je premica z enačbo $y=-\frac{4}{3}x+2$, ki seka koordinatni osi v točkah $M$ in $N$.

    \begin{parts}
        \part[4]
        Narišite premico.

            \begin{solution}[8cm]
                ???
            \end{solution}

        \part[4]
        Določite implicitno obliko enačbe vzporednice k dani premici, ki poteka skozi točko $(-2,1)$. % Premico tudi narišite.

            \begin{solution}[7cm]
                ???
            \end{solution}

        \part[3]
        Izračunajte razpolovišče daljice, ki jo določata točki $M$ in $N$.

            \begin{solution}[2cm]
                ???
            \end{solution}

    \end{parts}
    
    \newpage


    % 5. vprašanje
    \question
    Dana je funkcija $$f(x)=\begin{cases}
    -x-2, & x\leq 1 \\ 3x-6, & x>2
    \end{cases}.$$ 
    
    \begin{parts}
        \part[4]
        Narišite dano funkcijo.

            \begin{solution}[8cm]
                ???
            \end{solution}

        \part[3]
        Izračunajte ničle funkcije in določite območje, kjer je funkcija naraščajoča.

            \begin{solution}[5cm]
                ???
            \end{solution}

    \end{parts}
    
    


        
    \question[4]
    Narišite dano množico točk v ravnini.
    % \begin{parts}
    %     \part[4]
        $$[-2,\infty)\times \{-1,3,4\}$$

            \begin{solution}[2cm]
                ????
            \end{solution}




\end{questions}



\end{document}